%% 第9巻　波と変化を読みとく　本文（第15回編集委員会）
%% 章は \bkinputchap{ch9-NN} で読む（まだない章は、組むときに飛ばして「BKPAGE missing」を出す）
%% 第9巻は巻の色がむらさきなので、部扉の絵の線は白（bkOn）でかく
\providecommand\bkinputchap[1]{\directlua{collectgarbage("collect") texio.write_nl("log","BKMEM #1 "..math.floor(collectgarbage("count")/1024).." MB")}\IfFileExists{#1.tex}{\input{#1}}{\clearpage\typeout{BKPAGE missing #1}}}
\bkinputchap{ch9-00}

%% ======================== 第1部 フーリエ解析 ========================
\bkPartOpener{1}{第9巻\hspace{0.8em}第1部}{フーリエ解析}%
  {172/{sin をたしていくと、四角い波ができる？},
   160/{ほしい成分だけを取り出すには？},
   148/{くり返さない音も、周波数に分けられる？}}%
  {% 方形波と、それに近づく sin の和（白い線）
   \begin{scope}[bkOn, line width=1.1pt, line cap=round, line join=round]
     \draw[line width=0.6pt] (116,126) -- (172,126);
     \draw[line width=0.8pt, dash pattern=on 2pt off 2pt] (118,126) -- (118,138) -- (130.5,138) -- (130.5,114) -- (143,114) -- (143,138) -- (155.5,138) -- (155.5,114) -- (168,114) -- (168,126);
     \draw[line width=1.6pt, domain=0:50, samples=200] plot ({118+\x}, {126+12*(sin(\x*14.4)+sin(3*\x*14.4)/3+sin(5*\x*14.4)/5)*1.08});
   \end{scope}}%
  {波を sin と cos の和に分けるのが、フーリエ級数とフーリエ変換だよ。かけて平均すると、ほしい成分だけが取り出せる理由と、音・電気・熱への使い方を見ていくよ。}
\bkPartContents{1}{%% 自動生成（gen.py）
\bktocchap{9-1}{波をたし合わせて形を作る：フーリエ級数}{sinをたしていくと、四角い波ができる}{4}{4}{4}{38}{C}
\bktocchap{9-2}{成分の取り出し方：直交性}{かけて積分すると、ほしい成分だけが残る}{4}{5}{5}{82}{C}
\bktocchap{9-3}{複素数で書くフーリエ級数}{sinとcosを、1つの式でまとめて書く}{4}{5}{5}{124}{C}
\bktocchap{9-4}{スペクトル：フーリエ変換}{くり返さない音も、周波数に分けられる}{7}{5}{5}{164}{C}
}%
{くり返す波をフーリエ級数で、くり返さない信号をフーリエ変換で、周波数の成分に分けられるようになります。熱方程式・波動方程式を、変数分離とフーリエ級数で解けるようになります。}
\bkpartbegin{1}{フーリエ解析}
\bkinputchap{ch9-01}
\bkinputchap{ch9-02}
\bkinputchap{ch9-03}
\bkinputchap{ch9-04}

%% ======================== 第2部 ラプラス変換 ========================
\bkPartOpener{2}{第9巻\hspace{0.8em}第2部}{ラプラス変換}%
  {172/{微分方程式が、ふつうの方程式に変わる？},
   160/{箱をつなぐと、どうなる？},
   148/{落ち着くかどうかは、どこを見ればわかる？}}%
  {% 2つの世界：t の世界のグラフ → s の世界の × 印（白い線）
   \begin{scope}[bkOn, line width=1.1pt, line cap=round, line join=round]
     \draw[line width=0.6pt] (116,112) -- (140,112);
     \draw[line width=0.6pt] (116,112) -- (116,140);
     \draw[line width=1.6pt, domain=0:22, samples=80] plot ({116+\x}, {112+24*exp(-\x/7)*cos(\x*40)});
     \draw[-{Stealth[length=2.8mm]}, line width=1.6pt] (142,126) -- (150,126);
     \draw[line width=0.6pt] (152,126) -- (174,126);
     \draw[line width=0.6pt] (166,112) -- (166,140);
     \foreach \y in {133,119} {\draw[line width=1.4pt] ({158-2},{\y-2}) -- ({158+2},{\y+2}); \draw[line width=1.4pt] ({158-2},{\y+2}) -- ({158+2},{\y-2});}
   \end{scope}}%
  {微分を $s$ のかけ算に変える「翻訳機」がラプラス変換だよ。変換表と部分分数で微分方程式を解き、伝達関数と極で、機械や回路のふるまいを見通すよ。}
\bkPartContents{2}{%% 自動生成（gen.py）
\bktocchap{9-5}{微分をかけ算に変える翻訳機：ラプラス変換}{微分方程式が、ふつうの方程式に変わる}{3}{5}{5}{214}{C}
\bktocchap{9-6}{変換表と性質}{よく使う関数の「翻訳表」を作ろう}{3}{5}{5}{248}{C}
\bktocchap{9-7}{ラプラス変換で微分方程式を解く}{部分分数に分けて、表で元にもどす}{3}{5}{5}{292}{C}
\bktocchap{9-8}{伝達関数：入れると出てくるを s で表す}{箱をつなぐと、伝達関数のかけ算になる}{3}{5}{5}{346}{C}
\bktocchap{9-9}{極と安定性：s の値でふるまいがわかる}{極の場所を見るだけで、落ち着くかどうかがわかる}{3}{5}{5}{388}{C}
}%
{ラプラス変換を、定積分・変換表・部分分数を使い分けて計算し、回路やばね・制御の微分方程式を解けるようになります。伝達関数と極の位置から、落ち着くかどうかを読めるようになります。}
\bkpartbegin{2}{ラプラス変換}
\bkinputchap{ch9-05}
\bkinputchap{ch9-06}
\bkinputchap{ch9-07}
\bkinputchap{ch9-08}
\bkinputchap{ch9-09}

%% ======================== 第3部 行列と微分方程式 ========================
\bkPartOpener{3}{第9巻\hspace{0.8em}第3部}{行列と微分方程式}%
  {172/{たくさんの微分方程式を、1行で書ける？},
   160/{行列の「指数関数」って、何？},
   148/{音楽の圧縮やエアコンの制御に、どう使われている？}}%
  {% 倒立振子（白い線）
   \begin{scope}[bkOn, line width=1.1pt, line cap=round, line join=round]
     \draw[line width=0.6pt] (118,110) -- (170,110);
     \fill (137,111.6) circle (1.6); \fill (151,111.6) circle (1.6);
     \draw[line width=1.4pt] (132,113.6) rectangle (156,120.4);
     \draw[line width=1.8pt] (144,120.4) -- (150,143.5);
     \fill (150,143.5) circle (2.4);
     \draw[-{Stealth[length=2.6mm]}, line width=1.1pt] (128,117) -- (120,117);
     \draw[-{Stealth[length=2.6mm]}, line width=1.1pt] (160,117) -- (168,117);
   \end{scope}}%
  {たくさんの量が同時に変わるしくみを、行列1つの式（状態方程式）で書くよ。行列の指数関数と固有値で、倒立振子が立っていられるかどうかまでわかるよ。}
\bkPartContents{3}{%% 自動生成（gen.py）
\bktocchap{9-10}{連立微分方程式を行列で：状態方程式}{たくさんの微分方程式を、1行で書く}{5}{5}{5}{444}{C}
\bktocchap{9-11}{行列の指数関数と固有値}{固有値が、落ち着くかどうかを決める}{5}{5}{5}{492}{C}
\bktocchap{9-12}{フーリエとラプラスの使いみち}{音楽の圧縮、機械の診断、エアコンの制御}{7}{3}{5}{538}{R}
\bktocchap{9-13}{工学の数学の地図}{7本の路線は、どこで合流した？}{all}{3}{5}{556}{E}
}%
{連立微分方程式を状態方程式で書き、行列の指数関数と固有値で解いて、安定かどうかを判定できるようになります。フーリエとラプラスの使いみちを、本物の機械で見ます。}
\bkpartbegin{3}{行列と微分方程式}
\bkinputchap{ch9-10}
\bkinputchap{ch9-11}
\bkinputchap{ch9-12}

%% ======================== 終章 ========================
\bkinputchap{ch9-13}

%% ======================== 巻末 ========================
\bkinputchap{ch9-back}
